% ==============================================================================
% 03_CHAPTER3_PART2_TESTING.TEX - KIỂM THỬ HỆ THỐNG VÀ ĐÁNH GIÁ KẾT QUẢ
% ==============================================================================

\section{Kiểm thử hệ thống (System Testing)}
\label{sec:kiem_thu_he_thong}

\subsection{Mục tiêu và phương pháp kiểm thử}
Kiểm thử phần mềm là giai đoạn then chốt nhằm phát hiện sớm các khiếm khuyết kỹ thuật, thẩm định tính chính xác của các quy tắc nghiệp vụ và đo lường độ tin cậy của hệ thống trước khi đưa vào vận hành thực tế tại điểm bán. Các mục tiêu kiểm thử bao gồm:
\begin{enumerate}[label=\arabic*., leftmargin=1.5cm]
    \item \textbf{Thẩm định tính đúng đắn về chức năng (Functional Verification):} Đảm bảo mọi tính năng từ đăng nhập, bán hàng, tính thuế VAT, đánh giá khuyến mãi đến thanh toán trực tuyến đều hoạt động đúng theo các ca sử dụng đã đặc tả tại Chương 2.
    \item \textbf{Đảm bảo tính toàn vẹn tài chính và kế toán:} Kiểm tra nghiêm ngặt việc cấm xóa đơn hàng đã hoàn tất (\texttt{COMPLETED}), kiểm tra tính chính xác của giao dịch bù trừ hoàn kho (\texttt{IN}) khi hủy đơn chờ.
    \item \textbf{Kiểm tra tính an toàn thông tin (Security Verification):} Kiểm định cơ chế kiểm soát truy cập theo vai trò (RBAC), khả năng chống tấn công phát lại mã phiên (Replay Attack) của cơ chế Token Rotation, và tính toàn vẹn của chữ ký số Webhook PayOS.
    \item \textbf{Đo lường hiệu năng phản hồi (Performance Verification):} Xác nhận thời gian phản hồi tại quầy thu ngân đạt tiêu chuẩn dưới $1$ giây đối với mọi thao tác bán hàng.
\end{enumerate}

Phương pháp kiểm thử chủ đạo được áp dụng là \textbf{Kiểm thử Hộp đen (Black-box Testing)} kết hợp cùng \textbf{Kiểm thử Tích hợp (Integration Testing)} đa tầng giữa ứng dụng ReactJS SPA, máy chủ Spring Boot API và các dịch vụ đám mây bên thứ ba (PayOS, AWS S3).

\subsection{Môi trường và bộ dữ liệu kiểm thử}
Hệ thống được kiểm thử trên môi trường phần cứng và dữ liệu chuẩn:
\begin{itemize}[leftmargin=1.5cm]
    \item \textbf{Thiết bị thử nghiệm máy khách (Client):} Máy tính xách tay CPU Intel Core i7, 16GB RAM, kết nối mạng LAN/Wi-Fi 5GHz; màn hình độ phân giải Full HD ($1920\times 1080$), trình duyệt Google Chrome phiên bản 128. Kết nối máy in nhiệt Xprinter chuẩn 80mm qua cổng USB.
    \item \textbf{Môi trường máy chủ (Server):} Hệ điều hành Windows 11 / Docker Desktop Linux Container, máy chủ ảo hóa Java 25 Runtime, MySQL Server 8.0 với InnoDB Engine.
    \item \textbf{Bộ dữ liệu thử nghiệm chuẩn (Test Dataset):}
    \begin{itemize}
        \item 5 danh mục hàng hóa (Cà phê, Trà sữa, Nước ép, Bánh ngọt, Đồ ăn vặt).
        \item 30 sản phẩm cơ sở với 65 biến thể SKU (phân loại Size S/M/L, Màu sắc).
        \item 4 nhóm tùy chọn Modifier (Topping, Mức đường, Mức đá, Kem Cheese).
        \item 5 chương trình khuyến mãi thử nghiệm (1 Happy Hour khung giờ vàng 14h-17h, 1 BOGO mua Trà sữa size L tặng size M, 2 mã Coupon giảm 20\% và 50.000đ).
        \item 15 hồ sơ khách hàng thân thiết CRM với các hạn mức chi tiêu tích lũy khác nhau.
        \item 3 tài khoản người dùng: 1 Quản trị viên (\texttt{admin@billing.com}) và 2 Thu ngân (\texttt{staff1@billing.com}, \texttt{staff2@billing.com}).
    \end{itemize}
\end{itemize}

\subsection{Ma trận kịch bản kiểm thử chi tiết (Test Cases Matrix)}
Dưới đây là bảng ma trận 15 ca kiểm thử đại diện cho các luồng nghiệp vụ cốt lõi và các tình huống ngoại lệ của hệ thống POS (Bảng \ref{tab:test_cases_matrix}).

\begin{longtable}{|c|p{3.2cm}|p{3.8cm}|p{4.2cm}|c|}
\caption{Ma trận kịch bản kiểm thử chức năng và an ninh hệ thống} \label{tab:test_cases_matrix} \\
\hline
\textbf{Mã ca} & \textbf{Tên ca kiểm thử} & \textbf{Dữ liệu \& Các bước thực hiện} & \textbf{Kết quả mong đợi} & \textbf{Kết quả} \\ \hline
\endfirsthead
\hline
\textbf{Mã ca} & \textbf{Tên ca kiểm thử} & \textbf{Dữ liệu \& Các bước thực hiện} & \textbf{Kết quả mong đợi} & \textbf{Kết quả} \\ \hline
\endhead
\hline
\endfoot
\hline
\endlastfoot

\textbf{TC01} & Đăng nhập với mật khẩu không chính xác & Nhập email: \texttt{staff1@billing.com}, mật khẩu: \texttt{WrongPass123!}, nhấn ''Đăng nhập''. & Trả mã HTTP 401. Giao diện hiển thị toast lỗi tiếng Việt: ''Email hoặc mật khẩu không chính xác''. Không cấp token. & \textbf{PASS} \\ \hline

\textbf{TC02} & Đăng nhập thành công với vai trò Thu ngân & Nhập email: \texttt{staff1@billing.com}, mật khẩu đúng, nhấn ''Đăng nhập''. & Trả mã HTTP 200. Access Token lưu trong RAM. Cookie \texttt{refreshToken} (HttpOnly) được tạo. Tự động điều hướng vào \texttt{/explore}. & \textbf{PASS} \\ \hline

\textbf{TC03} & Kiểm soát phân quyền RBAC (\texttt{AdminRoute}) & Tài khoản Thu ngân cố tình truy cập thủ công vào URL quản trị: \texttt{/users} hoặc \texttt{/promotions}. & Hệ thống phát hiện vai trò không phải \texttt{ROLE\_ADMIN}, hiển thị toast cảnh báo từ chối truy cập và tự động đẩy về \texttt{/dashboard}. & \textbf{PASS} \\ \hline

\textbf{TC04} & Tìm kiếm sản phẩm thông minh (Debounce) & Tại thanh tìm kiếm quầy POS, gõ chuỗi ký tự ''tra sua''. & Sau khoảng nghỉ $300$ms, danh sách sản phẩm lọc đúng các mặt hàng trà sữa. Tốc độ lọc mượt mà, không giật lag. & \textbf{PASS} \\ \hline

\textbf{TC05} & Chọn Biến thể SKU và tùy chọn Modifier & Chọn Trà sữa Oolong, mở modal: chọn Size L (+$5.000$đ), chọn Trân châu trắng (+$10.000$đ), chọn 50\% đường. Nhấn thêm giỏ. & Mặt hàng được đưa vào giỏ với đơn giá đã cộng chính xác phụ thu ($55.000$đ). Hiển thị chi tiết lựa chọn kèm theo. & \textbf{PASS} \\ \hline

\textbf{TC06} & Nhận diện khách hàng thân thiết CRM & Nhập số điện thoại \texttt{0987654321} vào ô tìm kiếm khách hàng tại giỏ hàng. & Hệ thống tìm kiếm theo thời gian thực, hiển thị đúng họ tên khách hàng, số điểm tích lũy và tổng chi tiêu trước đó. & \textbf{PASS} \\ \hline

\textbf{TC07} & Tự động áp dụng khuyến mãi Happy Hour & Đặt đơn hàng trị giá $120.000$đ trong khung giờ 14h00 - 17h00 (chương trình giảm 15\% cho đơn từ $100$k). & Hệ thống tự động tính chiết khấu $18.000$đ. Tiền hàng chịu thuế giảm còn $102.000$đ. Thuế VAT 10\% là $10.200$đ. Tổng tiền là $112.200$đ. & \textbf{PASS} \\ \hline

\textbf{TC08} & Áp dụng mã Coupon hợp lệ và hết hạn & 1. Nhập mã \texttt{SALE20} hợp lệ $\rightarrow$ được giảm 20\%. \newline 2. Nhập mã \texttt{HETHAN50} đã hết hạn $\rightarrow$ nhấn áp dụng. & 1. Áp dụng mã thành công, hiển thị chiết khấu. \newline 2. Báo lỗi: ''Mã giảm giá đã hết hạn sử dụng'', giỏ hàng giữ nguyên giá trị gốc. & \textbf{PASS} \\ \hline

\textbf{TC09} & Thanh toán Tiền mặt và In hóa đơn nhiệt 80mm & Tổng đơn $112.200$đ. Thu ngân chọn Tiền mặt, bấm nút gợi ý nhận $200.000$đ từ khách. Nhấn ''Hoàn tất''. & Hệ thống tính tiền thừa thối lại: $87.800$đ. Mở popup in bill nhiệt 80mm với đầy đủ thông tin, kích hoạt lệnh in trình duyệt. & \textbf{PASS} \\ \hline

\textbf{TC10} & Trừ kho sổ cái khi tạo đơn hàng mới & Đơn hàng bán 2 cốc Trà sữa Oolong Size L (kho hiện có 50). Hoàn tất thanh toán. & Bảng sổ cái ghi nhận 1 giao dịch \texttt{OUT} với số lượng là $-2$. Trường \texttt{cachedStockQuantity} của biến thể giảm xuống 48. & \textbf{PASS} \\ \hline

\textbf{TC11} & Thanh toán trực tuyến VietQR PayOS & Chọn phương thức \texttt{PAYOS}. Mở popup QR. Khách quét mã chuyển tiền qua App Ngân hàng. & PayOS gửi Webhook xác thực SHA-256 Checksum hợp lệ. Đơn chuyển \texttt{COMPLETED}. Màn hình POS tự động nhận và mở bill. & \textbf{PASS} \\ \hline

\textbf{TC12} & Chuyển đổi phương thức thanh toán (\textit{Switch to Cash}) & Khách chọn PayOS nhưng app ngân hàng bảo trì. Thu ngân nhấn ''Chuyển sang tiền mặt''. & Đơn hàng chuyển sang \texttt{CASH} \texttt{COMPLETED}. Link PayOS bị hủy từ xa. Hệ thống hoàn tất đơn mà không nhân đôi xuất kho. & \textbf{PASS} \\ \hline

\textbf{TC13} & Hủy đơn hàng PENDING và Bù trừ kho tự động & Đơn hàng chờ chuyển khoản bị khách hủy. Thu ngân nhấn nút ''Hủy đơn hàng'' tại màn hình Lịch sử đơn. & Đơn chuyển \texttt{CANCELLED}. Sổ cái tự động tạo giao dịch \texttt{IN} (+2) hoàn trả kho. Hoàn lại 1 lượt dùng voucher. Giảm trừ CRM. & \textbf{PASS} \\ \hline

\textbf{TC14} & Chống xóa đơn hàng đã hoàn tất (\texttt{COMPLETED}) & Thu ngân hoặc hacker cố tình gửi request \texttt{DELETE /api/v1.0/orders/\{id\}} với đơn hàng đã thanh toán. & Máy chủ từ chối với mã lỗi 400 Bad Request: ''Nghiêm cấm xóa đơn hàng đã hoàn tất''. Bảo vệ nguyên vẹn dữ liệu tài chính. & \textbf{PASS} \\ \hline

\textbf{TC15} & Ngăn chặn Replay Attack với Refresh Token (ADR-0007) & Kẻ gian đánh cắp và cố tình sử dụng lại một Refresh Token cũ đã từng được xoay vòng trước đó. & Máy chủ phát hiện token đã thu hồi (\texttt{revokedAt} khác null), lập tức vô hiệu hóa toàn bộ nhóm \texttt{familyId}, buộc đăng nhập lại. & \textbf{PASS} \\
\end{longtable}

\subsection{Đánh giá và tổng kết kết quả kiểm thử}
Kết quả thực nghiệm trên toàn bộ 15 kịch bản kiểm thử trọng yếu (Bảng \ref{tab:test_summary}) khẳng định:
\begin{itemize}[leftmargin=1.5cm]
    \item \textbf{Độ tin cậy chức năng đạt 100\%:} Toàn bộ các ca kiểm thử nghiệp vụ bán hàng, quản lý biến thể, tính toán tài chính và tích hợp cổng thanh toán VietQR PayOS đều đạt kết quả mong đợi (15/15 Pass).
    \item \textbf{Hiệu năng đáp ứng vượt trội:} Thời gian xử lý trung bình của máy chủ API Spring Boot cho các tác vụ tạo đơn hàng và trừ kho sổ cái chỉ dao động trong khoảng $120$ms – $250$ms; thời gian cập nhật giao diện Front-end đạt dưới $30$ms, đáp ứng hoàn hảo tiêu chí phản hồi dưới $1$ giây đặt ra ban đầu.
    \item \textbf{Tính toàn vẹn kế toán và an toàn dữ liệu được bảo vệ tuyệt đối:} Quy trình bù trừ giao dịch kho khi hủy đơn (\texttt{compensatePendingOrder}) và cơ chế ngăn chặn xóa đơn đã hoàn tất đảm bảo số liệu doanh thu và tồn kho của cửa hàng luôn minh bạch và chính xác tuyệt đối.
\end{itemize}

\begin{table}[htbp]
\centering
\caption{Bảng tổng hợp kết quả kiểm thử hệ thống}
\label{tab:test_summary}
\vspace{2mm}
\begin{tabularx}{\textwidth}{|p{4.5cm}|c|c|c|X|}
\hline
\textbf{Nhóm kiểm thử} & \textbf{Tổng số ca} & \textbf{Đạt (Pass)} & \textbf{Không đạt} & \textbf{Tỷ lệ thành công} \\ \hline
Xác thực, Phân quyền \& Bảo mật & 4 & 4 & 0 & 100\% \\ \hline
Nghiệp vụ Bán hàng POS \& CRM & 4 & 4 & 0 & 100\% \\ \hline
Tính toán Khuyến mãi \& Thuế VAT & 2 & 2 & 0 & 100\% \\ \hline
Thanh toán PayOS \& Chuyển đổi & 2 & 2 & 0 & 100\% \\ \hline
Quản lý Kho sổ cái \& Toàn vẹn & 3 & 3 & 0 & 100\% \\ \hline
\textbf{TỔNG CỘNG} & \textbf{15} & \textbf{15} & \textbf{0} & \textbf{100\% (ĐẠT CHUẨN)} \\ \hline
\end{tabularx}
\end{table}

% ==============================================================================
\section{Tổng kết chương 3}
Chương 3 đã hoàn thành xuất sắc việc trình bày quá trình cài đặt, hiện thực hóa và kiểm thử thực nghiệm toàn bộ hệ thống POS Quản lý bán hàng. Các giải pháp kiến trúc đã được hiện thực hóa sống động qua 11 màn hình giao diện thực tế của cả phân hệ Thu ngân và Quản trị, từ giao diện bán hàng nhanh, popup biến thể SKU, thanh toán VietQR động đến mẫu in nhiệt 80mm. Năm thuật toán cốt lõi về quản lý kho theo sổ cái, tính toán tài chính phía server, bộ lọc JWT, đánh giá khuyến mãi no-stacking và xác thực chữ ký số Webhook PayOS đã được phân tích minh bạch qua các đoạn mã nguồn Java thực tế. Cuối cùng, ma trận 15 kịch bản kiểm thử toàn diện đã chứng minh tính ổn định, độ bảo mật cao và hiệu năng phản hồi nhanh chóng của hệ thống, khẳng định tính khả thi và giá trị ứng dụng thực tiễn cao của đề tài.
